%! Modeling with Quadratics — Unit 3, Lesson 3.7 — Group Activity (Answer Key)
\documentclass[10pt]{article}
\usepackage{algebra2-article}
\usepackage{algebra2-key}   % pulls in -boxes; do NOT also load -boxes
\newcommand{\TallMath}[1]{$\displaystyle #1\rule[-1.4em]{0pt}{3.2em}$}

\begin{document}

\pageheader{Unit 3, Lesson 3.7}{Group Activity --- Answer Key}
\namepartnerperiod

\vspace{0.08in}

\begin{headlinebox}{forestbg}
\small\textbf{Model It!} With your partner, keep the strategy in view: \textbf{define} the variable, \textbf{write}
the quadratic, decide \textbf{which feature} the question wants (start $\to$ $y$-intercept, ``reaches $0$'' $\to$
zero, ``the most/best'' $\to$ vertex), \textbf{find} it, and \textbf{interpret} it with units. Throw out any answer
that can't happen in the story.
\end{headlinebox}

\vspace{0.06in}

\begin{tcolorbox}[colback=white, colframe=black!40, title=\textbf{Tier R --- Remediate},
    fonttitle=\bfseries, arc=1.5mm, left=3mm, right=3mm, top=2mm, bottom=2mm, breakable]
\textbf{1. Read the projectile graph.} A ball is kicked straight up from the ground. Use the graph of
$h(t)=-16t^2+64t$ to answer.
\begin{center}
\begin{tikzpicture}[xscale=1.05, yscale=0.05]
  \draw[->, semithick, charcoal] (-0.2,0)--(4.6,0) node[below right, font=\scriptsize]{$t$ (s)};
  \draw[->, semithick, charcoal] (0,-3)--(0,74) node[left, font=\scriptsize]{$h$ (ft)};
  \draw[very thick, forest, smooth, samples=60, domain=0:4] plot (\x,{-16*\x*\x+64*\x});
  \draw[dashed, linegray] (2,0)--(2,64)--(0,64);
  \node[circle, fill=charcoal, inner sep=1.4pt] at (0,0){};
  \node[circle, fill=charcoal, inner sep=1.4pt] at (2,64){};
  \node[circle, fill=charcoal, inner sep=1.4pt] at (4,0){};
  \node[font=\scriptsize, anchor=west] at (2.08,64){vertex $(2,64)$};
\end{tikzpicture}
\end{center}
\begin{itemize}\setlength{\itemsep}{3pt}
  \item Starting height ($y$-intercept): \ans{$0$ ft (kicked from the ground)}
  \item Maximum height (vertex output): \ans{$64$ ft}
  \item Time it reaches the maximum (vertex input): \ans{$2$ s}
  \item When it returns to the ground (positive zero): \ans{$4$ s}
\end{itemize}

\textbf{2. Match the question to the feature.} Draw a line from each question to the feature that answers it.
\begin{center}
\renewcommand{\arraystretch}{1.3}
\begin{tabularx}{0.92\linewidth}{@{}X c X@{}}
``How high does it start?''        & $\rightarrow$ & the \textbf{$y$-intercept} \\
``What's the greatest value?''     & $\rightarrow$ & the \textbf{vertex} \\
``When does it hit the ground?''   & $\rightarrow$ & a \textbf{zero} (positive root) \\
\end{tabularx}
\end{center}
\ansline{Start $\to$ $y$-intercept; greatest value $\to$ vertex; hits the ground $\to$ zero.}

\textbf{3. Interpret a given area model.} A pen's area is $A(x)=-2x^2+40x$; its vertex is $(10,200)$. Fill in:
the maximum area is \ans{$200$} ft$^2$, using a width of \ans{$10$} ft.
\end{tcolorbox}

\begin{tcolorbox}[colback=white, colframe=black!40, title=\textbf{Tier A --- Approaching Proficiency},
    fonttitle=\bfseries, arc=1.5mm, left=3mm, right=3mm, top=2mm, bottom=2mm, breakable]
\textbf{1. Projectile.} A diver leaps upward at $32$ ft/s from a $128$-ft cliff: $h(t)=-16t^2+32t+128$.
\[ \text{max height} = \ans{$144$ ft} \qquad \text{time of max} = \ans{$1$ s} \qquad \text{hits water at } t=\ans{$4$ s} \]

\textbf{2. Maximum area (three sides).} A garden along a wall uses $80$ ft of fence on three sides ($x=$ width).
Write $A(x)$, then find the width and the maximum area.
\[ A(x)=\ans{$-2x^2+80x$} \qquad x=\ans{$20$ ft} \qquad \text{max area}=\ans{$800$ ft$^2$ ($20\times40$)} \]

\textbf{3. Maximum revenue.} A shop sells $300$ shirts at \$12. Each \$1 increase loses $15$ sales
($x=$ number of \$1 increases). Write $R(x)$, then find the best price and the maximum revenue.
\[ R(x)=\ans{$-15x^2+120x+3600$} \qquad \text{price}=\ans{$\$16$} \qquad \text{max }R=\ans{$\$3840$} \]
\end{tcolorbox}

\begin{tcolorbox}[colback=white, colframe=black!40, title=\textbf{Tier E --- Extension},
    fonttitle=\bfseries, arc=1.5mm, left=3mm, right=3mm, top=2mm, bottom=2mm, breakable]
\textbf{1. Optimize with a domain.} A rectangular garden against a wall uses $36$ ft of fence on three sides. Find the
width that maximizes the area, the maximum area, \emph{and} the feasible domain of $A(x)$.
\[ x=\ans{$9$ ft} \qquad \text{max area}=\ans{$162$ ft$^2$} \qquad \text{feasible domain}=\ans{$0<x<18$} \]

\textbf{2. Same height (systems tie-back).} Object A: $h_A(t)=-16t^2+48t$. Object B rises steadily: $h_B(t)=16t$.
Other than launch, when are they at the \emph{same} height, and what height is that? \emph{Hint: set the expressions equal.}
\[ t=\ans{$2$ s} \qquad \text{height}=\ans{$32$ ft} \]

\textbf{3. Interpret.} For the diver in Tier A \#1, explain what $h(0)=128$ means in context, and why the vertex's
\emph{output} (not its input) is the maximum height.
\ansline{$h(0)=128$ is the diver's starting height --- the $128$-ft cliff. The vertex is $(1,144)$: its input $t=1$
is \emph{when} the peak happens, and its output $144$ is \emph{how high} --- so the maximum height is the output, $144$ ft.}
\end{tcolorbox}

\begin{teachernote}
Tier R stays graphical/verbal: reading start/max/landing off the curve, matching questions to features, and reading a
\emph{given} vertex --- no algebra. Tier A is the build-and-solve engine; insist on units and on the vertex output vs.\
input distinction (\#1: max is $144$ ft, \emph{at} $t=1$). \#2's frequent error is $A=x(80-x)$ --- the barn wall means
$80-2x$. \#3 sets revenue as price\,$\times$\,\emph{changing} demand ($R=-15x^2+120x+3600$, vertex $x=4$, price \$16).
Tier E \#1 adds the feasible domain $0<x<18$ (an honest ``interpret in context'' beat); \#2 reuses substitution from
Lesson 3.6 --- the two objects meet at $t=2$, both at $32$ ft. Early finishers: ask them for Object A's \emph{own}
maximum height ($t=1.5$, $36$ ft) and whether the objects are ever at the same height while A is falling.
\end{teachernote}

\end{document}
